\documentclass[12pt]{article}
\usepackage[utf8]{inputenc}
\usepackage[T1]{fontenc}
\usepackage{hyperref}
\title{Long Wat Language Resources for K-12 Curriculum Development}
\author{Compiled by Claude Sonnet 5 (claude-sonnet-5) for Andr\'as Kornai}
\date{}
\begin{document}
\maketitle

\section{Introduction}
Long Wat (ISO 639-3 ttw; Glottocode long1406) is a Kenyah language of Sarawak, Malaysia. It is also called Tutoh and is spoken along the Tutoh River and in the village of Long Wat. Wikipedia reports about 600 native speakers, citing a 1981 figure, and says many speakers are shifting to Malay. Do not confuse it with Kiput, which is sometimes called Long Tutoh but is a separate language. Very little was found for this language.

\section{Primary Educational Resources}
No primers, textbooks or school materials were found.

\section{Linguistic and Reference Materials}
\begin{itemize}
\item Tutoh language (Wikipedia): \url{https://en.wikipedia.org/wiki/Tutoh_language}
\item Glottolog entry for the Nuclear Long Wat node (the Glottocode shown there is nucl1560): \url{https://glottolog.org/resource/languoid/id/nucl1560.bigmap.html}
\item Austronesian Comparative Dictionary, "Kenyah (Long Wat)" entries drawn from Blust (1971): \url{https://acd.clld.org/values/19298-157c868836217ecc164cd65c8c6a0d6d-1}
\end{itemize}

\section{Religious and Community Resources}
Nothing relevant was found.

\section{Digital and Audio Resources}
\begin{itemize}
\item Endangered Languages Project entry (it appeared in the results, but I did not confirm that it is the right language): \url{https://endangeredlanguages.com/lang/9769?hl=de}
\end{itemize}

\section{Curriculum Development Context}
No evidence of school use was found. Any curriculum work would need new documentation, probably in cooperation with the community and related Kenyah language work. The resources found here are minimal.
\end{document}